\documentclass[10pt]{article}

\def\Type{LessonDJ}
\def\TypeNum{Workshop Week 10} %Lesson #
\def\TypeTitle{Optimization and Net Change } %Lesson Title
\def\Semester{Fall 2021} %Not Used 
\def\Course{MATH 1190}%Course number and name
\def\Targets{  \target{A2}, \target{I1} }

\usepackage{preambleDJ} 

%%%%%%%%%%%% BEGIN DOCUMENT %%%%%%%%%%%%%%%
%%%%%%%%%%%% BEGIN DOCUMENT %%%%%%%%%%%%%%%
%%%%%%%%%%%% BEGIN DOCUMENT %%%%%%%%%%%%%%%
%%%%%%%%%%%% BEGIN DOCUMENT %%%%%%%%%%%%%%%
%%%%%%%%%%%% BEGIN DOCUMENT %%%%%%%%%%%%%%%

\begin{document}
\pagestyle{otherpages}
\thispagestyle{firstpage}


\vs{.6}
\
%\begin{center}
%    \includegraphics[width=0.8\textwidth]{IMG_2891.jpeg}
%\end{center}


\begin{enumerate}

\item {\bf Optimization} \\
\mpageT{.5}{.6}{
An athletic field is shaped like a rectangle with a semicircular track on each of its shorter sides (see Figure 1).  
The total perimeter of the track is fixed at $400$ meters.   The variable $R$ represents the radius  of each semi-circle.  The variable $L$ is the length  of the straight portion of the track.\\\\
Find the dimensions that enclose the largest possible area.
}{ \centering
\begin{tikzpicture}[scale=1.2, thick]
  % Parametri
  \def\R{1.5} % raggio dei semicerchi
  \def\L{4}   % lunghezza del rettangolo
%
  % Parte rettangolare (linee orizzontali)
  \draw (0,\R) -- (\L,\R);
  \draw (0,-\R) -- (\L,-\R);
  \draw[dotted] (0,\R) -- (0,0);
  \draw[dotted] (\L,\R) -- (\L,0);
  \node[circle, fill=black, inner sep=1pt] at (0,0) {};
  \node[circle, fill=black, inner sep=1pt] at (\L,0) {};
%
%
  % Semicerchi alle estremità
  \draw (0,-\R) arc[start angle=270, end angle=90, radius=\R];
  \draw (\L,\R) arc[start angle=90, end angle=-90, radius=\R];
%
  % Etichette
  \node at (\L/2, \R-0.4) {$L$};
  \node at (-\R+1.3, \R/2) {$R$};
  \node at (\L+\R-1.3, \R/2) {$R$};
  \node at (\L/2, -\R-0.4) {};
  %\node at (-0.3, \R+0.8) {Figure 1: Track shape};
%
\end{tikzpicture}
\vs{-.3}\captionof{figure}{Track Shape}
%\captionof{figure}{Track Shape}
}

{\bf Modeling Phase}
\tasksA{2}{
\task Write an expression for the total area enclosed by the track in terms of $L$ and $R$. Recall that the area of a half-circle is $A=\pi R^2/2$
\task Write an expression for the total perimeter in terms of your variables, and set it equal to $400$. Recall that the perimeter of a half-circle in $P=\pi R$\vs{.5}
\task Use the perimeter constraint to express one variable in terms of the other, and substitute it into the area formula.
}



\vs{.5}

{\bf Calculus Phase}
 Think about the process you have been using to optimize other quantities. Use that process here to find the dimensions that maximize the area of the athletic field.

\newpage
\item {\bf Net Change and The Definite Integral} .\\
Consider the graph below, which represents the velocity (in mi/hr) of a braking vehicle (looks like one of those ``rolling stops") over time (in seconds s). 
%In addition, suppose this function is given by $f(x) = 10\log\left(2\pi\ x\right).$
\graphicsC{bus_go_broooom.png}{4}

%Consider the following questions:
    \begin{enumerate}
    \item What does the definite integral $\ds \int_5^{15} f(x) dx$ represent? In particular, what do the bounds symbolize \textbf{ and } how would you explain what the integral calculates (give units).\\
        \vs{1.5}
    \item Using at least two rectangles, what is the approximate area under this curve from $x = 5$ to $x=15$? What does this quantity represent (use units)?\\
        \vfill

%    \item Using at least one rectangle, approximate 
%        \begin{enumerate}
%            \item $\ds \int_5^{10} f(x) dx$
%                \vfill
%            \item $\ds \int_{10}^{15} f(x) dx$
%                \vfill
%            \item $\ds \int_{15}^{20} f(x) dx$
%                \vfill
%        \end{enumerate}
%    \item What can you say about the difference (using subtraction) between (i.) and (ii.). What about the difference between (ii.) and (iii.)? Which difference is greater? How do you suspect $\ds \int_{20}^{25}f(x) dx$ might relate to (iii.)?
%                \vfill
    \end{enumerate}

\newpage
\item {\bf Optimization} The sum of two numbers is $50$. Find the two numbers that maximizes the product. Let's use , $x$ and $y$ to represent our two numbers.
%\[
%f(x,y) = xy
%\]
%subject to the constraint
%\[
%x^2 + 4y^2 = 16.
%\]

{\bf Modeling Phase}
\tasksA{2}{
\task Write down a formula for the product, $P$ of our two numbers ($x$ and $y$)
\task Write down a formula that represents ``The sum of two numbers is $50$." This should involve both $x$ and $y$.  This formula is called a {\em constraint}.
\task Use the constraint to express one variable in terms of the other.
\task Substitute your expression for the product, $P$.  \( f(x,y) \) to obtain a single-variable function.
\task Simplify \( P \) and determine the domain of possible \(x\)-values.
}

\vs{1}

{\bf Calculus Phase}
\\ Now that you have obtained $P$ as a single-variable function, use the usual techniques to find the maximum and the minimum of $P$ (if any). 


\vs{1}


\newpage 

\mpageT{.5}{.5}{
\item {\bf Optimization} Suppose Bob wants to build a fence to enclose some sheep with a rectangular sheep pen. Fortunately, he lives next to a river, so he only needs to build three sides of the enclosure. Bob has $30$m of fencing. What are the dimensions that will enclose the greatest amount of sheep?
}{\graphicsC{sheepPen}{3}
}
{\bf Modeling Phase}
\tasksA{2}{
\task Create variables that describe the width and length of the sheep pen.
\task What quantity are we going to maximize? Let's call this quantity $S$.  Use both of your new variables to write down an equation for this quantity?\vs{1}
\task Bob's constraint is that he only has $30$m of fencing. Use your variables to write down an equation that represents this constraint.
\task Both the quantity you are going to maximize, $S$, and your constraint contain both of your variables. Use your constraint to solve for one variable and substitute  into $S$.
}
\vs{1}
{\bf Calculus Phase} Think about the process you have been using to optimize other quantities. Use that process here to find the dimension that maximize the area of the sheep pen?


\newpage
\item {\bf Optimization} Suppose John is standing next to a river that is 1 mile wide. His house is on the opposite of the river 3 miles down. John can swim at 5 miles per hour but he can run at 8 miles per hour. What is the shortest amount of time it can take John to reach his house? You will find this by finding the best choice for $d$ where John stops running and starts swimming.\\
\mpageT{.4}{.6}[\hs{.1}]{
{\bf Modeling Phase} We first need to develop a model of the total time: $$T_{total}=T_{run}+T_{swim}$$
The total horizontal distance to the house is $3$ miles.  Once John starts to swim, the remaining horizontal distance is $d$ miles (see Figure 1). 
}{\graphicsC{runSwim}{4}{Figure 1}
}
\tasksA{2}{
\task What is the distance John ran? Your answer should contain $d$.
\task How long does it take for John to run this distance? This is $T_{run}$.  Your answer should contain $d$.\vs{.5}
\task What is the distance John swam?  {\em Hint}: You will want to utilize Pythagorean's Theorem to find this distance. Your answer should contain $d$.
\task How long does it take for John to swim this distance? This is $T_{swim}$.  Your answer should contain $d$. \vs{.75}
\task* What is $T_{total}$? Your answer will contain $d$\vs{.25}
}
{\bf Calculus Phase} Think about the process you have been using to optimize quantities. Use that process to find the minimum time it takes for John to get home. You will definitely want to use technology for the last part of finding critical numbers.


    \vfill



\newpage











\end{enumerate}

\end{document}